\unnumberedchap{Introduction}
\noindent\rule{\textwidth}{0.4pt}

D'après la \textit{Délégation générale à la langue française et aux langues de France} (DGLFLF) : 
{« }certaines langues [de France] ne disposent pas des ressources et outils permettant les interactions 
courantes avec les usagers telles que la traduction automatique, la synthèse vocale, la génération 
automatique de texte, ou encore la rédaction assistée.{ »}\footnote{
    \hyperlink{https://www.culture.gouv.fr/thematiques/langue-francaise-et-langues-de-france/agir-pour-les-langues/innover-dans-le-domaine-des-langues-et-du-numerique/soutenir-et-encourager-la-diversite-linguistique-dans-le-domaine-numerique}
    {Soutenir et encourager la diversité linguistique (culture.gouv.fr)}
}
Dans ce contexte où les langues régionales de France disposent de très peu de ressources numériques 
exploitables. La constitution de corpus parallèles représente un défi majeur. Or, ces corpus sont 
essentiels pour l’entraînement de systèmes de traduction automatique, d’alignement ou de représentation 
sémantique multilingue.

Face à la rareté des corpus parallèles existants pour des langues comme le corse, l’occitan ou l’alsacien, 
une question centrale viens alors à se poser : comment identifier automatiquement des fragments parallèles exploitables 
à partir de corpus comparables non alignés ?

L’identification de tels fragments constitue une étape intermédiaire stratégique. Elle permet d’enrichir progressivement 
des ressources parallèles, d’améliorer les modèles de traduction neuronale et de soutenir la représentation numérique de 
langues faiblement dotées.

Ce mémoire vise donc à présenter, de la manière la plus complète possible, l'état de l'art et 
la réalisation d'un programme informatique répondant aux attentes de notre sujet. 
Nous allons commencer par expliquer les problématiques liées aux langues peu dotées et certaines
des solutions déjà apportées pour pallier ces problèmes. Ensuite, nous expliquerons les fondamentaux des études précédemment
réalisées sur les corpus comparables. Nous présenterons aussi certaines des réalisations passées pour
les méthodes d'alignement. Enfin, nous discuterons des réalisations liées au mémoire en expliquant le fond
théorique ainsi que la réalisation du projet principal.